\Article{32}
How the magnetic fluids are actually distributed within the Earth necessarily remains undetermined. Indeed, by a general theorem already mentioned in Article 2 of the \textit{Intensitas}, which is to be treated in detail on another occasion, any distribution of the magnetic fluids within a three-dimensional region can always be replaced by a distribution on its surface in such a way that the effect at every point outside the region remains exactly the same. It follows readily that \textit{one and the same} effect throughout the exterior can arise from infinitely many \textit{different} distributions of the magnetic fluids in the interior.

We can, however, specify the fictitious distribution on the Earth's surface that is entirely equivalent to the actual distribution in the interior with respect to the forces it produces outside the Earth. Owing to the Earth's \SourceLeaf{170}spherical shape, we can even do so in a very simple way. The density of the magnetic fluid at every point on the Earth's surface. that is, the amount of fluid per unit area, is expressed by the formula
\[
\frac1{4\pi}(\frac VR-2Z)
\]
or by
\[
-\frac1{4\pi}(3P'+5P''+7P'''+9P''''+\text{etc.})
\]
The value of this formula will be illustrated in due course by a graph. Here it need only be noted that the value is negative in the northern half of the Earth and positive in the southern half, with the dividing line crossing the equator twice (at $6^0$ and $186^0$ longitude) and reaching as far as about $15^0$ north and south latitude on either side. There are \textit{two} minima in the northern half and only one maximum in the southern half. A rough calculation gives the following values and positions for these minima and the maximum
\begin{center}
\begin{tabular}{rlrlrl}
$-209{,}1$&at&$55^0$&N. latitude&$263^0$&longitude\\
$-200{,}0$&at&$71^0$&N. latitude&$116^0$&longitude\\
$+277{,}7$&at&$70^0$&S. latitude&$154^0$&longitude
\end{tabular}
\end{center}
The values themselves are based on the unit used for our elements. They must therefore still be multiplied by $0{,}0034941$ if they are to be expressed in absolute measure.

\Article{33}
As already stated above, our elements are to be regarded solely as a first approximation. As such, they agree sufficiently well with the observations, as Article 29 shows. There is no doubt that refining the calculation on the basis of these observations would bring about much closer agreement. The only difficulty with such a calculation would be its length, which remains daunting even if one tries to shorten it by methods like those astronomers use to refine the orbital elements of planets and comets. Although this difficulty could readily be
\SourceLeaf{171}
overcome if the work could be divided among a number of people doing the calculations, it may be inadvisable to make such a refinement now, when the data from so many places that would have to be included are still so unreliable. For the time being it will be best to continue comparing the elements with observations. This will provide a way of making the general maps much more reliable than was possible with the exclusively empirical method used so far. May we be allowed, however, to take a few glimpses at the theory's future progress, whose full realization may still be a long way off.

\Article{34}
To refine and complete the elements satisfactorily, the observational data must meet much higher requirements than have been met so far. At every point used, the data should have a precision thus far attained at only a very few points. They should be cleared of irregular movements and should apply to the same instant. It will be a long time before these requirements can be met. What is most urgently needed now is to obtain \textit{complete} observations (that is, observations covering all three elements) at one point or another within those large regions where such observations are still entirely lacking. Indeed, the farther a new point lies from the others already at our disposal, the more important it is for the general theory.

After a sufficient interval, the elements will be determined afresh for a second epoch, yielding their secular changes. For this purpose it will clearly be indispensable to abandon the measure of intensities used so far altogether and replace it with an absolute measure.

Over the coming centuries, these changes too will cease to appear uniform, and investigating how the elements change will furnish natural scientists with an inexhaustible subject for research.

\SourceLeaf{172}
\Article{35}
The future will also shed light on interesting points of the theory.

Our theory assumes that each measurable magnetized part of the body of the Earth contains exactly as much positive as negative fluid. If the magnetic fluids had no reality and were merely a fictitious substitute for galvanic currents in the smallest parts of the Earth, the validity of that substitution would itself entail this equality. If, however, we attribute real existence to the magnetic fluids, we could reasonably question whether the amounts of the two fluids are exactly equal. For individual magnetic bodies (natural or artificial magnets), very precise experiments could easily settle whether they contain a perceptible excess of one fluid or the other. If there were such an excess, a plumb line weighted with such a body would have to deviate from the vertical (specifically, in the direction of the magnetic meridian). If experiments of this kind, carried out with many artificial magnets at a site sufficiently far from iron, were never to show the slightest deviation (as seems likely), this equality could indeed be assumed for the whole Earth with the greatest probability, though the possibility of some inequality would still be open.

The only difference such an inequality would make to our theory is that $P^0$ (Article 17) would no longer be $=0$ The consequence would be that throughout the infinite exterior the expression for $Z$ would require the additional term $\frac{RRP^0}{rr}$, and thus the (constant) term $P^0$ would have to be added at the Earth's surface, while $X$ and $Y$ would be entirely unaffected. When the future has provided a much larger stock of precise observations than is now available, it will indeed be possible to determine whether their accurate representation requires a nonzero value of $P^0$ or not. Given the present state of the data, however, such an undertaking could still yield no result at all.
